\documentclass{test-template}

\begin{document}

\begin{center}
    \textbf{\Huge{Experimental Mathematics}} \\
    Exercises Sheet \#1 \\
    Error-Correction Codes
\end{center}

\subsection*{Motivation}

One of the most useful cryptographic gadget that frequently arises in numerous
applications is \textbf{Error-Correction Codes} (ECC). While we further give the
rigorous definition, the rough purpose of ECC is, as the name suggests, correct
errors during transmission through noisy channel or during some complicated
processing of data. As previously mentioned, such construction is used in
various applications: ranging from error-correction in computer memory or data
storage devices (such as HDDs, SSDs, USBs) to fuzzy extractors (recovering
strong cryptographic keys from biometrics) or zero-knowledge protocols (namely,
zk-STARKs).

Error-Correction Codes are primarily based on linear algebra and polynomials.
The primary goal of this exercise sheet is to implement the \textit{Reed-Solomon
Codes}, which are among most basic and fundamental types of ECC. 

\subsection*{Theory}

We define the $(n,k)$ error-correction code $\mathcal{C}$ first. Fix the finite
alphabet $\Sigma$. You might think of alphabet as the basic building blocks
which constitute messages and codewords. Most commonly $\Sigma$ is just
$\{0,1\}$ (in such case, both messages and codewords are nothing but a sequence
of bits) or, for linear codes, some finite field elements $\mathbb{F}_q$ (more
on that later).

We call error-correction code $\mathcal{C}$ an $(n,k)$-code if messages come
from set $\Sigma^k$, while codewords from $\Sigma^n$. Always assume $n \geq k$:
we always extend $k$-sized messages to the larger $n$-sized codespace. This way,
if some errors occur in larger-dimensional space, we hope that we can fix a
certain number of errors.

Now we want to induce certain metric over $\Sigma^n$. This is done using 
\textbf{Hamming distance}. Let us give the definition.

\begin{definition}
    The \textbf{Hamming distance} $\Delta(\mathbf{x},\mathbf{y})$ for two
    codewords $\mathbf{x},\mathbf{y} \in \Sigma^n$ is defined as the number of
    positions where $\mathbf{x}$ and $\mathbf{y}$ differ:
    \begin{equation*}
        \Delta(\mathbf{x},\mathbf{y}) \triangleq \#\{i \in [n]: x_i \neq y_i\}, \quad \mathbf{x},\mathbf{y} \in \Sigma^n.
    \end{equation*}

    The Hamming distance naturally induces the norm over $\Sigma^n$ denoted by
    $\|\mathbf{x}\|_0 := \Delta(\mathbf{x},\mathbf{0})$, which is called a
    \textbf{Hamming weight}.
\end{definition}

\end{document}